\paperchapter{July 2006}{July 2006}{A}{50}{GSET 2006}
\keytable{\kcell{1}{B} & \kcell{2}{C} & \kcell{3}{A} & \kcell{4}{A} & \kcell{5}{C} & \kcell{6}{B} & \kcell{7}{A} & \kcell{8}{B} & \kcell{9}{A} & \kcell{10}{D} \\
\kcell{11}{D} & \kcell{12}{B} & \kcell{13}{C} & \kcell{14}{A} & \kcell{15}{A} & \kcell{16}{D} & \kcell{17}{A} & \kcell{18}{C} & \kcell{19}{D} & \kcell{20}{B} \\
\kcell{21}{A} & \kcell{22}{C} & \kcell{23}{C} & \kcell{24}{C} & \kcell{25}{D} & \kcell{26}{D} & \kcell{27}{D} & \kcell{28}{B} & \kcell{29}{A} & \kcell{30}{D} \\
\kcell{31}{C} & \kcell{32}{B} & \kcell{33}{D} & \kcell{34}{A} & \kcell{35}{B} & \kcell{36}{C} & \kcell{37}{A} & \kcell{38}{B} & \kcell{39}{A} & \kcell{40}{A} \\
\kcell{41}{C} & \kcell{42}{B} & \kcell{43}{A} & \kcell{44}{C} & \kcell{45}{D} & \kcell{46}{D} & \kcell{47}{B} & \kcell{48}{D} & \kcell{49}{A} & \kcell{50}{B}}
\begin{question}{1}
{Which one of the following is the most important characteristics of a good teacher?\par}
\optstwo{To control the class effectively.}{To motivate the students for learning}{To have command over the subject}{To use different teaching methods.}
\answer{B}{To motivate the students for learning}
\begin{explain}
The most important quality of a good teacher is the ability to \textbf{motivate students to learn}. Motivation drives attention, effort and lifelong learning; mastery of subject and methods serve this end.
\end{explain}
\end{question}

\begin{question}{2}
{During the teaching-learning process a teacher should\par}
\optsone{provide relevant information}{maintain a lively classroom environment}{allow students to raise questions}{give relevant illustrations.}
\answer{C}{allow students to raise questions}
\begin{explain}
During teaching, a teacher should \textbf{allow students to raise questions}. Questions show engagement, clarify doubts and make learning two-way.
\end{explain}
\end{question}

\begin{question}{3}
{The ultimate goal of education is to\par}
\optsone{develop problem-solving skills among students}{make students economically productive}{cultivate the habit of maintaining good relationships}{prepare knowledgeable individuals.}
\answer{A}{develop problem-solving skills among students}
\begin{explain}
Per the key, the ultimate goal of education is to \textbf{develop problem-solving skills} so that learners can face real-life situations independently.
\end{explain}
\end{question}

\begin{question}{4}
{The main function of teachers in higher education is\par}
\optstwo{teaching, research and extension.}{being a friend, philosopher and guide.}{to solve problems of students.}{to generate knowledge.}
\answer{A}{teaching, research and extension.}
\begin{explain}
The main functions of teachers in higher education are \textbf{teaching, research and extension} (service to the community), the three pillars of universities.
\end{explain}
\end{question}

\begin{question}{5}
{If a student asks you a question in the class the answer of which you do not know, then what will you do?\par}
\optsone{Get angry with the student.}{Do not allow anyone to ask questions.}{Accept your lack of knowledge and tell that you will come prepared with the same next time.}{Will feel guilty.}
\answer{C}{Accept your lack of knowledge and tell that you will come prepared with the same next time.}
\begin{explain}
An honest teacher \textbf{accepts the lack of knowledge} and promises to come prepared next time, modelling integrity and a learning attitude.
\end{explain}
\end{question}

\begin{question}{6}
{Which one of the following is not the aim of scientific research?\par}
\optstwo{Theory building.}{Logical reasoning.}{Problem solving.}{Creating new knowledge.}
\answer{B}{Logical reasoning.}
\begin{explain}
Scientific research aims at theory building, problem solving and creating new knowledge. \textbf{Logical reasoning} is a \emph{means} (a tool) of research, not its aim.
\end{explain}
\end{question}

\begin{question}{7}
{An investigator studied some selected subjects for certain years continuously. The type of research is:\par}
\optstwo{Longitudinal study.}{Cross-sectional study.}{Experimental study.}{Correlational study.}
\answer{A}{Longitudinal study.}
\begin{explain}
Studying the same subjects repeatedly over several years is a \textbf{longitudinal study}. A cross-sectional study observes different groups at one point in time.
\end{explain}
\end{question}

\begin{question}{8}
{Which one is the most important factor for selection of a research problem?\par}
\optstwo{Researcher's interest.}{Researchability of the problem.}{Administrative suitability.}{Uniqueness of the problem.}
\answer{B}{Researchability of the problem.}
\begin{explain}
The most important factor is \textbf{researchability}: whether the problem can be studied empirically with available data, tools, time and resources.
\end{explain}
\end{question}

\begin{question}{9}
{Inductive reasoning implies\par}
\optstwo{particular to general.}{general to particular.}{symbolic logic.}{none of the above.}
\answer{A}{particular to general.}
\begin{explain}
\textbf{Inductive} reasoning moves from \textbf{particular to general} (specific observations $\rightarrow$ general conclusion). Deduction moves from general to particular.
\end{explain}
\end{question}

\begin{question}{10}
{Which one of the following is not the example of qualitative research?\par}
\optstwo{Case study.}{Ethnography.}{Content analysis.}{Pretest-Posttest experiment.}
\answer{D}{Pretest-Posttest experiment.}
\begin{explain}
Case study, ethnography and content analysis are qualitative. A \textbf{pretest--posttest experiment} is a quantitative experimental design.
\end{explain}
\end{question}

\begin{context}
\textbf{Direction (Q.~No.~11 to 15):} Read the following text and answer questions 11 to 15.\par\smallskip
Sixty per cent of rural India lives in utter poverty with no electricity. Rural India uses 180 million tones of fuel every year for cooking, creating indoor pollution and health hazards. The WHO estimates that 1.5 million deaths per year are caused by pollution from chulha smoke. Making available locally grown clean and renewable cooking and lighting fuel can improve the quality of rural life. Liquid fuels are far superior to solid fuels for cooking because of their clean burning and higher energy. Ethanol is one of the best as it is an excellent substitute for kerosene and burns better with no unpleasant smell. Its burning power is almost as clean as that of LPG.

Ethanol can be produced from any sugary material and is presently produced world over from sugarcane and corn. However, as the ethanol economy grows, there is need to produce it from a crop which uses much less water than sugarcane and also produces food. Sorghum (Jowar) is one such crop whose stem is sweet like sugarcane. Its earhead produces grains which can be used for making bread. Its sweet stem has nearly the same amount of sugar and hence the juice can be fermented and used for ethanol production. The left-over stem after juice extraction, together with leaves is excellent fodder for animals. So from the same price of land one can get food, fuel and fodder. Besides, sweet sorghum (Jowar) uses nearly fifty percent less water than sugarcane to produce the same amount of sugar. It is a four month crop so farmers can grow two crops per year from the same piece of land. Also the energy output ratio from sweet sorghum is very positive.
\end{context}
\begin{question}{11}
{The purpose of the paragraph is to bring out:\par}
\optstwo{Problems of rural poor.}{The limitations of sugar cane.}{Benefits of sorghum (Jowar).}{Remedies for pollutant energy.}
\answer{D}{Remedies for pollutant energy.}
\begin{explain}
The paragraph identifies a problem (polluting solid cooking fuels) and proposes a solution: clean, locally grown ethanol from sweet sorghum. Its purpose is to bring out \textbf{remedies for pollutant energy}.
\end{explain}
\end{question}

\begin{question}{12}
{Most suitable heading for the paragraph will be:\par}
\optstwo{Rural poverty.}{Alternate Fuels.}{Solid Fuels Vs. Liquid Fuels.}{Efficient Fuels.}
\answer{B}{Alternate Fuels.}
\begin{explain}
The passage argues for \textbf{alternate fuels} (ethanol from sweet sorghum) in place of polluting solid fuels.
\end{explain}
\end{question}

\begin{question}{13}
{As per the authors, who are the rural poor?\par}
\optsone{Who do not have any electricity.}{Who use solid fuel for cooking.}{Who live with pollution and health hazards.}{Who do not use cooking gas.}
\answer{C}{Who live with pollution and health hazards.}
\begin{explain}
Per the key, the rural poor are those \textbf{who live with pollution and health hazards} (from chulha smoke, with no electricity).
\end{explain}
\end{question}

\begin{question}{14}
{As per the author which can be easy source of energy\par}
\optsone{fuel that can be locally produced.}{fuel that can be produced from crops.}{fuel that can be used as alternate to kerosene.}{fuel that does not give unpleasant smell.}
\answer{A}{fuel that can be locally produced.}
\begin{explain}
``Making available \textbf{locally grown} clean and renewable cooking and lighting fuel can improve the quality of rural life.''
\end{explain}
\end{question}

\begin{question}{15}
{The nature of the paragraph is\par}
\optsfour{analytical.}{critical.}{conceptual.}{comprehensive.}
\answer{A}{analytical.}
\begin{explain}
The paragraph analyses the problem and compares crops and fuels using evidence (water use, energy output ratio), so it is \textbf{analytical}.
\end{explain}
\end{question}

\begin{question}{16}
{Students are inattentive in the classroom. The reason is:\par}
\optstwo{Lack of interest in the subject.}{Outside obstacles.}{Nature of the subject.}{Poor communication.}
\answer{D}{Poor communication.}
\begin{explain}
Inattention in class usually signals \textbf{poor communication} by the teacher: unclear, monotonous or irrelevant presentation.
\end{explain}
\end{question}

\begin{question}{17}
{Communication is a network of\par}
\optstwo{Interactions.}{Messages.}{Receivers and senders.}{Information.}
\answer{A}{Interactions.}
\begin{explain}
Communication is a network of \textbf{interactions} between people: a dynamic two-way process, not just messages or information.
\end{explain}
\end{question}

\begin{question}{18}
{Which one of the following is not a barrier to communication ?\par}
\optstwo{Noise.}{Culture.}{Job satisfaction.}{Attitude.}
\answer{C}{Job satisfaction.}
\begin{explain}
Noise, culture and attitude can be barriers. \textbf{Job satisfaction} is not a barrier to communication.
\end{explain}
\end{question}

\begin{question}{19}
{Which of the following is true with respect to a good communicator?\par}
\optsone{May not always be a good teacher.}{Has humourous expression.}{Has a wise range of reading interest.}{Has mastery over the medium of instruction.}
\answer{D}{Has mastery over the medium of instruction.}
\begin{explain}
A good communicator \textbf{has mastery over the medium of instruction}, i.e.\ command of language and channel.
\end{explain}
\end{question}

\begin{question}{20}
{Which of the following are the important characteristics of an effective teacher ?\newline (i) Knowledge of content.\newline (ii) Effective verbal communication skills.\newline (iii) Impressive personality.\newline (iv) Ability to command respect from students.\par}
\optsfour{(ii) and (iii)}{(i) and (ii)}{(iii) and (iv)}{(i) and (iv).}
\answer{B}{(i) and (ii)}
\begin{explain}
The most important characteristics of an effective teacher are \textbf{knowledge of content (i)} and \textbf{effective verbal communication (ii)}.
\end{explain}
\end{question}

\begin{question}{21}
{Find the odd number out: 8, 9, 64, 125, 216, 49\par}
\optsfour{125}{216}{64}{49}
\answer{A}{125}
\begin{explain}
The numbers alternate cube and square: $2^3=8$, $3^2=9$, $4^3=64$, $5^2=25$, $6^3=216$, $7^2=49$. \textbf{125} breaks the pattern (it should be 25).
\end{explain}
\end{question}

\begin{question}{22}
{Which is the missing number in the following series ? 16, 54, 128, \blank{}\par}
\optsfour{162}{256}{250}{416.}
\answer{C}{250}
\begin{explain}
$2\times2^3=16$, $2\times3^3=54$, $2\times4^3=128$, so next $2\times5^3=\mathbf{250}$.
\end{explain}
\end{question}

\begin{question}{23}
{Fill in the blank: B, G, K, \blank{}, P\par}
\optsfour{L}{M}{N}{J.}
\answer{C}{N}
\begin{explain}
B(2) $+5$ G(7) $+4$ K(11) $+3$ \textbf{N(14)} $+2$ P(16).
\end{explain}
\end{question}

\begin{question}{24}
{If CAT = ECV, then SUN = ?\par}
\optsfour{NSU}{VCE}{UWP}{TAC.}
\answer{C}{UWP}
\begin{explain}
Each letter moves two places forward: C$\rightarrow$E, A$\rightarrow$C, T$\rightarrow$V. SUN $\rightarrow$ \textbf{UWP}.
\end{explain}
\end{question}

\begin{question}{25}
{Find the missing number: 95, 87, 80, 74, 69, \blank{}\par}
\optsfour{49}{59}{48}{65.}
\answer{D}{65.}
\begin{explain}
The differences fall by one: $-8, -7, -6, -5, -4$. So $69-4=\mathbf{65}$.
\end{explain}
\end{question}

\begin{question}{26}
{January 23, 2005 was a Sunday. What will be the day on December 31, 2004?\par}
\optsfour{Monday.}{Sunday.}{Saturday.}{Friday.}
\answer{D}{Friday.}
\begin{explain}
31 December 2004 is 23 days before 23 January 2005. $23=3\times7+2$, so it falls two days before Sunday: \textbf{Friday}.
\end{explain}
\end{question}

\begin{question}{27}
{Who is the odd man out? George Bush, Gen. Mushrraf, APJ Abdul Kalam, Tony Blair.\par}
\optstwo{George Bush}{Gen. Mushrraf}{APJ Abdul Kalam}{Tony Blair.}
\answer{D}{Tony Blair.}
\begin{explain}
George Bush (USA), Gen.\ Musharraf (Pakistan) and APJ Abdul Kalam (India) were then \textbf{heads of state (Presidents)}. \textbf{Tony Blair} was UK \emph{Prime Minister} (head of government), so he is the odd one out.
\end{explain}
\end{question}

\begin{question}{28}
{Shyam remembers that their mother visited them after 15th but before 20th of a month. Shyam's brother remembers that their mother had visited them after 18th but before 22nd. On which date did their mother visit them?\par}
\optsfour{16th}{19th}{21st}{17th.}
\answer{B}{19th}
\begin{explain}
Shyam: 16--19. Brother: 19--21. The common date is the \textbf{19th}.
\end{explain}
\end{question}

\begin{question}{29}
{If we substitute numbers 1 to 12 indicating hours on the dial of a clock by alphabets in the order starting with B, which letter will represent 9?\par}
\optsfour{J}{L}{K}{I.}
\answer{A}{J}
\begin{explain}
1$\rightarrow$B, 2$\rightarrow$C, \ldots, so $n\rightarrow$ the $(n+1)$th letter. 9 $\rightarrow$ 10th letter $=$ \textbf{J}.
\end{explain}
\end{question}

\begin{question}{30}
{The individual weights of 'A' and 'B' are the same as that of 'C'. 'C's weight is equal to 'D's weight. Whose weight is the maximum ?\par}
\optstwo{D}{C}{A}{All have the same weight.}
\answer{D}{All have the same weight.}
\begin{explain}
A $=$ C and B $=$ C (each individually equals C), and C $=$ D. So A $=$ B $=$ C $=$ D: \textbf{all have the same weight}.
\end{explain}
\end{question}

\begin{context}
Answer the questions 31 to 35 from the data given in the following table. In this table the runs made by the five batsmen in five matches are reported.
\begin{center}\small
\begin{tabular}{lccccc|c}\toprule \textbf{Batsman} & M1 & M2 & M3 & M4 & M5 & \textbf{Total}\\\midrule
Mohan & 44 & 53* & 26 & 22* & 38 & 183\\ Anand & 22 & 15 & 42 & 13 & 0 & 92\\ Kishor & 72 & 56 & 48 & 63 & 51 & 290\\ Chetan & 32 & 135 & 0 & 182 & 0 & 349\\ Ramesh & 09 & 11 & 17 & 06 & 10 & 53\\\bottomrule
\multicolumn{7}{l}{\scriptsize * Not out. (Totals column added for the solution.)}
\end{tabular}
\end{center}
\end{context}
\begin{question}{31}
{Who made highest total run ?\par}
\optstwo{Kishor}{Anand \& Ramesh}{Chetan}{Mohan \& Anand.}
\answer{C}{Chetan}
\begin{explain}
Totals: Mohan 183, Anand 92, Kishor 290, \textbf{Chetan 349}, Ramesh 53. Chetan made the highest total.
\end{explain}
\end{question}

\begin{question}{32}
{Who is consistently poor batsman?\par}
\optsfour{Anand}{Ramesh}{Mohan}{Chetan.}
\answer{B}{Ramesh}
\begin{explain}
\textbf{Ramesh} scored low in every match (9, 11, 17, 6, 10), so he is consistently poor.
\end{explain}
\end{question}

\begin{question}{33}
{Who made zero run more times ?\par}
\optstwo{Anand}{Ramesh}{Ramesh and Anand both}{Chetan}
\answer{D}{Chetan}
\begin{explain}
\textbf{Chetan} scored zero twice (Matches 3 and 5); Anand once.
\end{explain}
\end{question}

\begin{question}{34}
{Who is consistent good batsman?\par}
\optsfour{Kishor}{Mohan}{Chetan}{None of above.}
\answer{A}{Kishor}
\begin{explain}
\textbf{Kishor} scored 48--72 in every match: consistently good. Chetan's runs fluctuate wildly.
\end{explain}
\end{question}

\begin{question}{35}
{Average run of mohan is\par}
\optsfour{18.00}{36.60}{61.00}{69.80.}
\answer{B}{36.60}
\begin{explain}
Mohan's average (runs $\div$ innings, as per the key) $=183\div5=\mathbf{36.60}$.
\tip{In cricket, the batting average divides by \emph{dismissals} ($183\div3=61$). The key uses the simple arithmetic mean.}
\end{explain}
\end{question}

\begin{question}{36}
{The instrument which helps in receiving the web-page matter through internet is called\par}
\optsfour{Client}{Browser}{Server}{Power Point}
\answer{C}{Server}
\keynote{On the user's side, the software that receives and displays web pages is the \emph{browser}; the server is the machine that supplies them. The official key is (C).}
\begin{explain}
The official key gives \textbf{Server}: the web server is the computer that stores web pages and delivers (serves) their content over the Internet on request.
\end{explain}
\end{question}

\begin{question}{37}
{MS WORD is an example of\par}
\optstwo{Application Software}{System Software}{Operating Software}{Translating Programme.}
\answer{A}{Application Software}
\begin{explain}
MS Word is \textbf{application software}: it performs a user task (word processing).
\end{explain}
\end{question}

\begin{question}{38}
{The first webpage of any organization is called\par}
\optsfour{Portal}{Home Page}{Vortal}{Website.}
\answer{B}{Home Page}
\begin{explain}
The first page of a website is the \textbf{home page}.
\end{explain}
\end{question}

\begin{question}{39}
{The advantage of internet to a teacher is\par}
\optsone{He can update his knowledge and enrich it.}{He can counsel the rural people.}{He can mark attendance of students.}{He can evaluate his own knowledge.}
\answer{A}{He can update his knowledge and enrich it.}
\begin{explain}
The internet lets a teacher \textbf{update and enrich his or her knowledge} continuously.
\end{explain}
\end{question}

\begin{question}{40}
{Which of the following does not have any use in Local Area Network (LAN) ?\par}
\optsfour{Modem}{Printer}{CPU}{Monitor.}
\answer{A}{Modem}
\begin{explain}
A \textbf{modem} connects a computer to a WAN or the Internet via telephone lines. It is not needed inside a LAN, where printers, CPUs and monitors are used.
\end{explain}
\end{question}

\begin{question}{41}
{In which city did India win the Hockey World Cup in 1975?\par}
\optsfour{New Delhi}{Toronto}{Kuala Lumpur}{London.}
\answer{C}{Kuala Lumpur}
\begin{explain}
India won the Hockey World Cup in \textbf{1975} at \textbf{Kuala Lumpur}, Malaysia, beating Pakistan 2--1 in the final.
\end{explain}
\end{question}

\begin{question}{42}
{Which branch of the following is studied statistically through demography?\par}
\optstwo{Human society}{Human population}{Human life}{Human psychology.}
\answer{B}{Human population}
\begin{explain}
\textbf{Demography} is the statistical study of \textbf{human population}: size, structure, births, deaths and migration.
\end{explain}
\end{question}

\begin{question}{43}
{Which gas is useful as well as harmful to human being?\par}
\optsfour{Ozone}{Oxygen}{Carbon dioxide}{Methane.}
\answer{A}{Ozone}
\begin{explain}
\textbf{Ozone} is useful in the stratosphere (it blocks UV) but harmful at ground level (a pollutant causing respiratory problems).
\end{explain}
\end{question}

\begin{question}{44}
{In which year the department of environment was established in India?\par}
\optsfour{1970}{1972}{1982}{1986}
\answer{C}{1982}
\begin{explain}
The Department of Environment was created in 1980 and became the Ministry of Environment and Forests in 1985. Among the options, the key gives \textbf{1982}.
\end{explain}
\end{question}

\begin{question}{45}
{Which gas is responsible for Green House Effect ?\par}
\optstwo{Oxygen}{Ozone}{Carbon Monoxide}{Carbon Dioxide.}
\answer{D}{Carbon Dioxide.}
\begin{explain}
\textbf{Carbon dioxide} is the main greenhouse gas responsible for the enhanced greenhouse effect.
\end{explain}
\end{question}

\begin{question}{46}
{Which one of the following is a deemed university?\par}
\optstwo{Sardar Patel University, V.V. Nagar.}{M. S. University of Baroda}{Bhavnagar University, Bhavnagar.}{None of above.}
\answer{D}{None of above.}
\begin{explain}
Sardar Patel University, M.~S.~University of Baroda and Bhavnagar University are all \textbf{state universities}, so \textbf{none} is deemed.
\end{explain}
\end{question}

\begin{question}{47}
{Main feature of distance education is:\par}
\optstwo{Face to face learning}{Counselling sessions and study centres}{Compulsory attendance}{Easy to get the degree.}
\answer{B}{Counselling sessions and study centres}
\begin{explain}
The main feature of distance education is \textbf{counselling sessions and study centres}, with learning through self-study material and no compulsory attendance.
\end{explain}
\end{question}

\begin{question}{48}
{Who is the chair person of University Grants Commission?\par}
\optstwo{Yashpal}{Madhuriben Shah}{Arun Nigvekar}{Sukhdeo Thorat}
\answer{D}{Sukhdeo Thorat}
\begin{explain}
In 2006, \textbf{Prof.\ Sukhadeo Thorat} was UGC Chairman (2006--2011). Arun Nigavekar preceded him; Yash Pal was Chairman in 1986--91.
\end{explain}
\end{question}

\begin{question}{49}
{Which body governs and looks after technical education in India?\par}
\optsfour{AICTE}{NCTE}{MCI}{ICSSR}
\answer{A}{AICTE}
\begin{explain}
Technical education is governed by the \textbf{AICTE}.
\end{explain}
\end{question}

\begin{question}{50}
{Purpose of knowledge commission is to:\par}
\optsone{Find means to generate knowledge}{Find relevance of knowledge to society}{Find ways to generate income through knowledge.}{Find methods to create structure of knowledge.}
\answer{B}{Find relevance of knowledge to society}
\begin{explain}
The National Knowledge Commission (2005, chaired by Sam Pitroda) was set up to build excellence in the knowledge system for the 21st century. Per the key, its purpose is to \textbf{find the relevance of knowledge to society}.
\end{explain}
\end{question}
